\section{Related Work}

As shown in Table~\ref{tab:provenance-comparison}, existing provenance tools stop at file or dataset granularity. Chen et al.~\cite{chen2026finegrained} propose fine-grained sample-level traceability with contribution scores across ML pipelines but focus on training-stage data usage auditing rather than authorship attribution. OriginBlame coexists with tools like DVC and HuggingFace Datasets: DVC manages file versions, while ob manages record-level author provenance.

Machine unlearning research similarly assumes the forget set is known (Table~\ref{tab:unlearning-benchmarks}). Benchmarks TOFU~\cite{maini2024tofu} and OpenUnlearning use pre-defined forget sets, while MUSE~\cite{shi2025muse} evaluates unlearning on real-world corpora with known forget/retain splits. Unlearning methods like NPO, RMU, GradAscent, and exact unlearning algorithms~\cite{muresanu2025exact} directly receive forget sets as input. ForSId~\cite{dangelo2025forsid} (\S2) is the only work addressing forget-set identification, but operates post-hoc with model access. OriginBlame captures origin information at recording time, requiring no model access. \begin{table}[t]
	\centering
	\caption{Machine Unlearning Benchmarks: The Forget-Set Location Gap}
	\label{tab:unlearning-benchmarks}
	\resizebox{\columnwidth}{!}{%
		\begin{tabular}{lccc}
			\toprule
			\textbf{Benchmark/Method}              & \textbf{Forget Set Source}                  & \textbf{Identifies Location?} \\
			\midrule
			TOFU \cite{maini2024tofu}              & Pre-defined (fictitious authors)            & $\times$                      \\
			MUSE \cite{shi2025muse}                & Pre-defined (2 corpora)                     & $\times$                      \\
			OpenUnlearning                         & Pre-defined                                 & $\times$                      \\
			NPO/RMU                                & Pre-defined                                 & $\times$                      \\
			ForSId \cite{dangelo2025forsid}        & Post-hoc inference                          & Partial                       \\
			\rowcolor{blue!10}\textbf{OriginBlame} & \textbf{Author revocation $\to$ forget set} & \textbf{$\checkmark$}         \\
			\bottomrule
		\end{tabular}%
	}
\end{table}


This positioning complements methods like Attribute-to-Delete~\cite{georgiev2024attribute}, which use datamodels to simulate unlearning but still require the forget set as input.

IPFS's content-addressing design~\cite{trautwein2022ipfs} inspired OriginBlame's hash chain mechanism, but IPFS operates at chunk level (256\,KB blocks) while OriginBlame refines granularity to lines. OriginBlame uses plain JSONL rather than a special file system. Compared to git-blame, designed for code editing scenarios, OriginBlame targets generated data scenarios: supporting multi-source attribution where a single line may fuse contributions from multiple authors, and author-initiated revocation.
